Pour desserrer un écrou rouillé, il faut un moment de force d'au moins \qty{\v{tau0}}{N.m}. Un mécanicien tire sur l'extrémité d'une clé de \qty{\v{L}}{cm} (mesurée depuis le centre de l'écrou) avec une force de \qty{\v{F}}{N} qui fait un angle de \qty{\v{th}}{\degree} avec le manche de la clé.

\begin{enumerate}[label=\alph*)]
    \item Déterminez le moment de force exercé sur l'écrou. L'écrou se desserre-t-il ?
    \item Déterminez la force minimale à appliquer à l'extrémité de la clé, perpendiculairement au manche, pour desserrer l'écrou.
    \item Déterminez la force minimale nécessaire si le mécanicien garde l'angle de \qty{\v{th}}{\degree}.
\end{enumerate}
